\section*{Erd\H{o}s Problem 651: convex configurations in higher dimensions}

\paragraph{Statement.}
Let \(f_k(n)\) be the least number of points in general position in
\(\mathbb R^k\) that forces \(n\) points in convex position.
Must \(f_k(n)>(1+c_k)^n\) hold for some \(c_k>0\)?

\paragraph{Claim and source.}
For every fixed \(k\geq3\), the answer is no.
The external input is Pohoata and Zakharov,
\emph{Convex polytopes from fewer points},
\url{https://arxiv.org/abs/2208.04878}, Theorem 1.1:
for every \(\varepsilon>0\), all sufficiently large \(n\) satisfy
\[
 f_3(n)\leq\lceil2^{\varepsilon n}\rceil.
\]
We spell out the projection and quantifier arguments that give the
disproof in every higher dimension. The planar case is different;
no disproof for \(k=2\) is claimed.
See \url{https://www.erdosproblems.com/651}.

\paragraph{Projection preserves the needed hypotheses.}
Let \(X\subset\mathbb R^k\) be a finite general-position set of at least
\(k+1\) points, with \(k\geq3\).
Every four points of \(X\) are affinely independent: otherwise adding
another \(k-3\) points would yield \(k+1\) points in an affine subspace
of dimension at most \(k-1\), contrary to general position.

Choose a linear map \(L:\mathbb R^k\to\mathbb R^3\) for which every
four-point subset remains affinely independent. Such a map exists:
for each fixed four-point subset, the determinant of its three
image-difference vectors is a nonzero polynomial in the entries of
\(L\), since those original difference vectors are independent.
The zero sets of finitely many nonzero polynomials cannot cover
the whole space of linear maps. The resulting image is a set of
distinct points with no four coplanar.

If \(Y\subseteq X\) projects to points in convex position, then \(Y\)
itself is in convex position. Indeed, an expression
\(y=\sum_{z\in Y\setminus\{y\}}\lambda_z z\), with nonnegative
\(\lambda_z\) summing to one, would project to the same forbidden
convex combination for \(L(y)\). Thus, in the asymptotic range
relevant here,
\[
 f_k(n)\leq f_3(n).
\]

\paragraph{Excluding every positive exponential constant.}
Fix any proposed \(c_k>0\), and take
\(\varepsilon=\tfrac12\log_2(1+c_k)>0\).
For all sufficiently large \(n\),
\[
 f_k(n)\leq\lceil2^{\varepsilon n}\rceil
          <2^{2\varepsilon n}=(1+c_k)^n.
\]
This contradicts the proposed exponential lower bound, even if it
was required only eventually. Consequently no such positive
constant exists in dimension \(k\geq3\).

\paragraph{Scope.}
The disproof is complete relative to the stated Pohoata--Zakharov
theorem. The generic projection and the comparison of growth rates
are proved here; the deep three-dimensional theorem is not
re-proved or presented as a new result.

\paragraph{Completion estimate.}
\textbf{COMPLETION: 100\%}
